\chapter{Suporte Tecnológico}

\section{Plataforma de desenvolvimento móvel}
A escolha inicial foi pelo Flutter, por permitir uma base única de código para Android e iOS. O Flutter é um kit de ferramentas de código aberto criado pelo Google para o desenvolvimento de interfaces a partir de uma única base de código, usando a linguagem Dart \cite{flutter_docs}. Contudo, o processo de publicação para iOS exige um ambiente macOS para compilar, assinar e enviar o aplicativo, recurso de que não dispomos no momento. Buscando alternativas, adotamos o React Native em conjunto com o Expo. Além disso, utilizamos o EAS para construção e distribuição em nuvem, e adotamos TypeScript como linguagem principal.
\subsection{React Native}
O React Native é um framework de código aberto para desenvolvimento de aplicativos nativos a partir do ecossistema React. O React, por sua vez, é uma biblioteca JavaScript criada para construir interfaces de usuário por meio de \textit{componentes} reaproveitáveis, um modo de organizar a tela em partes menores e independentes \cite{react_docs}. O React Native aplica essas mesmas ideias no contexto móvel, permitindo criar aplicativos para Android e iOS usando JavaScript ou TypeScript, com acesso a componentes nativos do sistema e recarga rápida durante o desenvolvimento \cite{reactnative_docs}.
\subsection{Expo}
O Expo é uma plataforma e conjunto de SDKs\footnote{SDK significa \textit{Software Development Kit}, ou conjunto de bibliotecas e ferramentas que facilitam o uso de funcionalidades prontas em um aplicativo.} de código aberto que padroniza o ambiente do React Native e simplifica tarefas comuns do ciclo de desenvolvimento, como inicialização de projetos, acesso a APIs de dispositivo e testes em aparelhos físicos\cite{expo_docs}. Na prática, o Expo reduz a fricção de configuração e acelera a iteração de funcionalidades.
\subsection{Expo Application Services (EAS)}
O EAS é o conjunto de serviços em nuvem do Expo para construir, assinar, publicar e gerenciar aplicativos React Native. O EAS Build compila binários para Android e iOS na nuvem, o EAS Submit auxilia no envio às lojas\cite{eas_build,eas_suite}. Esse fluxo permite gerar builds para iOS sem a necessidade de um computador com macOS local.
\subsection{TypeScript}
O TypeScript é um superconjunto tipado de JavaScript que adiciona tipagem estática, inferência de tipos e recursos avançados de linguagem. No contexto do aplicativo, a tipagem melhora a legibilidade, ajuda a prevenir erros em tempo de desenvolvimento e favorece a manutenção em projetos de médio porte, com suporte nativo no ecossistema React e no Node.js \cite{typescript_handbook}.

\section{Backend e estilo de API}
O backend é a camada responsável por processar a lógica de negócio, manipular dados e disponibilizar serviços para o aplicativo móvel. Nesta etapa, optamos por tecnologias consolidadas que oferecem bom desempenho, comunidade ativa e ampla documentação.
\subsection{Node.js}
O Node.js é um ambiente de execução para JavaScript fora do navegador, construído sobre o motor V8 do Google Chrome \cite{node_docs}. Ele permite criar aplicações escaláveis e de alto desempenho utilizando a mesma linguagem tanto no cliente quanto no servidor. Além disso, conta com o ecossistema NPM\footnote{Node Package Manager, repositório público com milhares de bibliotecas reutilizáveis que aceleram o desenvolvimento.} que disponibiliza pacotes prontos para diversas finalidades, o que reduz tempo de implementação. A familiaridade da equipe com a tecnologia e a possibilidade de adotar TypeScript no servidor reforçaram essa escolha.
\subsection{Express}
O Express é um framework minimalista para aplicações web no Node.js \cite{express_docs}. Ele fornece recursos como roteamento de URLs e organização modular de controladores e serviços. Sua simplicidade, combinada com a flexibilidade, facilita a construção de APIs robustas sem adicionar complexidade desnecessária.
\subsection{REST}
A interface entre o backend e o aplicativo é feita por meio de uma API\footnote{Application Programming Interface. Em termos práticos, um conjunto de endpoints que define como clientes podem enviar requisições e receber respostas de um serviço.}. Adotamos o estilo arquitetural REST (Representational State Transfer), proposto por Roy Fielding em sua tese de doutorado \cite{fielding2000}. O REST não é um protocolo, mas um conjunto de princípios que orienta o design de sistemas distribuídos, normalmente implementados sobre HTTP \cite{rfc7231}. Utilizando métodos como \texttt{GET}, \texttt{POST}, \texttt{PUT} e \texttt{DELETE}, conseguimos padronizar a forma como os dados são enviados e recebidos. 
Entre suas vantagens estão a simplicidade de implementação, a compatibilidade com diferentes clientes e a escalabilidade\footnote{A escalabilidade é favorecida quando o serviço é \textit{stateless}, ou seja, não mantém estado entre requisições. Isso permite que diferentes servidores atendam às requisições de forma independente, facilitando balanceamento e réplica.}. Em resumo, o REST equilibra praticidade e robustez, atendendo às necessidades atuais do projeto com baixo atrito de integração e amplo suporte no mercado.

\section{Banco de dados}
% PostgreSQL na AWS Free Tier

\section{Serviços e SDKs externos}
% Google Maps SDK e Reverse Geocoding; deep link para Google Maps ou Waze

\section{Autenticação e notificações}
% Firebase Auth (e-mail e código); Expo Notifications para push

\section{Ferramentas de projeto e colaboração}

Para organizar, visualizar e documentar o trabalho utilizamos um conjunto de ferramentas que se complementam entre si. Cada uma cumpre uma função específica, e em conjunto elas reduzem o atrito no dia a dia, favorecem a comunicação e garantem o registro das principais decisões do projeto.

\subsection{Miro}
O Miro é um quadro branco digital colaborativo em tempo real \cite{miro_whiteboard}. Permite rascunhar ideias, construir fluxos, criar mapas visuais e conduzir dinâmicas de ideação de forma coletiva. Usamos como mural central para anotar insights surgidos em reuniões e esboçar protótipos de baixa fidelidade.

\subsection{Google Drive}
O Google Drive é um serviço de armazenamento em nuvem para documentos, planilhas e apresentações \cite{drive_workspace}. Utilizamos para guardar e compartilhar materiais do projeto, como questionários, identidade visual, análises e atas, garantindo acesso rápido e organização em pastas.

\subsection{Jira}
O Jira é uma plataforma de gestão ágil que disponibiliza quadros no estilo Kanban \cite{jira_intro}. Foi utilizada para manter o backlog atualizado e acompanhar o andamento das tarefas de desenvolvimento, oferecendo visibilidade ao time sobre o que está em progresso e o que já foi concluído.

\subsection{Figma}
O Figma é uma ferramenta online de design e prototipação de interfaces \cite{figma_site}. Com ela produzimos protótipos de alta fidelidade, registramos fluxos de navegação e validamos escolhas de interface de forma visual e interativa.

\subsection{GitHub}
O GitHub é uma plataforma de hospedagem de repositórios Git \cite{github_version_control}. Adotamos para versionar o código do aplicativo e do backend, permitindo controle de histórico, colaboração entre desenvolvedores e segurança dos artefatos de software.

Em conjunto, essas ferramentas sustentam o planejamento, a execução e a documentação do projeto em um espaço de trabalho unificado, acessível e colaborativo.


\section{Padrões de código e qualidade}
% ESLint e Prettier; builds com EAS; sem CI no TCC 1

\section{Segurança e LGPD no nível tecnológico}
% Confirmação de idade mínima (16+); aceite de Termos/Política; uso responsável de localização


